\documentclass[12pt]{article}
\usepackage[utf8]{inputenc}
\usepackage[T1]{fontenc}
\usepackage{hyperref}
\title{Nago, Northern Language Resources for K-12 Curriculum Development}
\author{Compiled by Claude Sonnet 5.5 (claude-sonnet-5-5) for Andr\'as Kornai}
\date{}
\begin{document}
\maketitle

\section{Introduction}
Northern Nago (ISO 639-3 xkb; Glottocode mani1294) is the Manigri-Kambol\'e variety of the Ede dialect continuum, spoken in Benin and Togo. Its ISO reference name was changed from ``Manigri-Kambol\'e Ede Nago'' to ``Northern Nago'' in 2010. Glottolog places it in Atlantic-Congo, Yoruboid, Edekiri, Ede. Joshua Project lists about 68,000 Manigri Nago in Benin and 73,000 Kambol\'e in Togo, though its page assigns the language to the Gur family, which conflicts with Glottolog. The search results show only catalogue entries and a small amount of audio or evangelistic material. They show no school materials, dictionary or grammar of this variety.

\section{Primary Educational Resources}
None found.

\section{Linguistic and Reference Materials}
\begin{itemize}
\item Glottolog languoid page for Manigri-Kambol\'e Ede Nago: \url{https://glottolog.org/resource/languoid/id/mani1294.bigmap.html}
\item Glottolog reference record that turned up in the search: \url{https://glottolog.org/resource/reference/id/144667}
\item ISO 639-3 resource record for xkb: \url{https://pure.mpg.de/cone/iso639-3/resource/xkb}
\item IANA language subtag template for xkb: \url{https://www.iana.org/assignments/lang-subtags-templates/xkb.txt}
\item Wikipedia, Ede language (the broader dialect continuum): \url{https://en.wikipedia.org/wiki/Ede_language}
\end{itemize}

\section{Religious and Community Resources}
\begin{itemize}
\item Joshua Project, Kambol\'e, Manigri Nago in Benin: \url{https://joshuaproject.net/index.php/people_groups/12472/BN}
\item Joshua Project language page for xkb: \url{https://joshuaproject.net/index.php/languages/xkb}
\end{itemize}
No scripture translation was confirmed.

\section{Digital and Audio Resources}
\begin{itemize}
\item Global Recordings Network, Kambole language page (xkb): \url{https://globalrecordings.net/en/language/xkb}
\end{itemize}
I did not check the contents of this page, so I cannot say what recordings it holds.

\section{Curriculum Development Context}
The language is stable and used by all in its community, but it is not known to be taught in schools. Benin and Togo each have a different official language (French) and their own national-language policies, which I did not research. The main gap is documentation: I found no grammar, dictionary, orthography or primer for this variety. Related Ede and Yoruba materials may eventually help but are separate languages or varieties and are not listed here.

\end{document}